\documentclass[11pt]{amsart}
\usepackage[a4paper,margin=27mm]{geometry}
\usepackage{amsmath,amssymb,amsthm,mathtools}
\usepackage{graphicx,booktabs,microtype,hyperref}
\usepackage{enumitem}
\usepackage{xcolor}
\hypersetup{colorlinks=true,linkcolor=blue,citecolor=blue,urlcolor=blue}
\newtheorem{conjecture}{Conjecture}
\newtheorem{proposition}{Proposition}
\newtheorem{remark}{Remark}
\DeclareMathOperator{\Nm}{N}
\newcommand{\PP}{\mathcal P}
\newcommand{\OO}{\mathbb Z[\omega]}
\newcommand{\rX}{R_X}
\newcommand{\gtilde}{\widetilde g}
\newcommand{\fig}[4]{%
 \begin{figure}[htbp]\centering
 \IfFileExists{figures/generated/#1.pdf}{\includegraphics[width=#2\linewidth]{figures/generated/#1.pdf}}{%
 \fbox{\parbox{.85\linewidth}{Run \texttt{python figures/make\_all.py} before compiling to generate figure \texttt{#1.pdf}.}}}
 \caption{#3}\label{fig:#4}\end{figure}}
\title[A shift-aspect crossover for sextic vectors]{A shift-aspect crossover for twisted sextic Gauss-sum vectors:\\
numerical evidence, conjectures, and a Lean-verified finite core}
\author{John Fairfax-Ball}
\date{27 October 2026}
\begin{document}
\begin{abstract}
We numerically investigate the action of the sextic residue-symbol matrix indexed
by primary split primes of the Eisenstein integers on vectors built from normalized
sextic Gauss sums and multiplicative twists. A dimensionless excess statistic
exhibits an empirical crossover in the ratio of sixth-power-free twist norm to
prime cutoff. Separately, fifth-power twists display a negative large-aspect
excess. The profiles were subjected to committed predictions, fresh twist-class
tests, prime-versus-squarefree-modulus controls, and independent cutoff tests up
to \(X=4\cdot10^6\); failed precursor predictions are retained. We state three
conjectures, carefully distinguishing compact-aspect convergence from unproved
fixed-class limits. A formally verified Lean 4 library establishes a finite
residue-field and Gauss-sum core, including a matrix-row identity; it does not
prove the numerical asymptotics. Exact signed Gram and reciprocity identities
identify the additional analytic input a proof would require.
\end{abstract}
\maketitle
\tableofcontents
\section{Introduction and results at a glance}\label{sec:intro}
For a cutoff \(X\), let \(\PP(X)\) index primary degree-one Eisenstein
primes and consider the complex matrix with entries
\((\rho/\sigma)_6\), extended by zero along the diagonal.
Multiplying this matrix by a Gauss-sum vector twisted by
\(\kappa\) produces an arithmetic analogue of a structured spectral
probe. We study the quotient between the output norm and the
random-phase scale. The central observation is a \emph{measured},
not proved, crossover as the dimensionless parameter
\(t=\Nm_6(\kappa)/X\) varies.

\fig{visual_abstract}{.96}{Schematic overview of the arithmetic
object and statistic. This panel is a diagram, not experimental
evidence; formulas use the conventions of CAND-002 in
\texttt{research/CONJECTURE\_REGISTER.md} at source commit
\texttt{30965467}.}{visual}
\fig{crossover}{.98}{The primary empirical profile. Exponents
\(1\)--\(4\) and the frozen RS-S016 discovery-bin means are displayed
against \(\log_{10}t\); transparent points are finite computed classes
from RS-S015, RS-S016 and RS-S021. Sources:
\texttt{results/rs\_s015\_class\_tests\_2026-10-26.json},
\texttt{rs\_s016\_x\_dependence\_2026-10-26.json} and
\texttt{rs\_s021\_large\_x\_2026-10-26.json}, at
\texttt{30965467}. No extrapolation beyond the measured points is made.}{collapse}

The candidate comprises three distinct \emph{conjectures}:
a compact-\(t\) crossover, a large-\(t\) fifth-power deficit,
and a separate fixed-class asymptotic. The third is \textbf{not}
a consequence of the first two. Small positive \(t\) and the boundary
\(t\to0\) are different regimes of quantification. The figure is not
a proof of any limiting statement.

The finite mathematical core consists of exact symbol and Gauss-sum
relations formalized in Lean, and paper-level identities reducing the
strength statistic to a signed two-point kernel. A proof of the observed
shape would require estimates for moving-shift prime-modulus first
moments and their \emph{signed} correlations, not a fit to a figure.

The work relates to residue-matrix characterizations
\cite{ddk2016}, metaplectic Gauss-sum theory
\cite{patterson1987,kazhdan1984,broeker2013,chinta2003},
the cubic and quartic first-moment problems
\cite{dunn2024,ddhl2023}, and recent results on large-sieve
nonorthogonality and sextic norms
\cite{dfdh2026,xusang2026}. These citations provide context, not
a claim of priority or an assertion that the present asymptotics
are established.

\section{Arithmetic setup and computational method}\label{sec:setup}
\subsection{Primes, residue symbols and normalisation}
Put \(\omega=e^{2\pi i/3}\) and \(\zeta_6=1+\omega\).
We use the repository's primary convention
\(\rho=a+b\omega\), \(a\equiv2\pmod3\), \(b\equiv0\pmod3\),
with rational prime norm \(p=\Nm(\rho)\le X\).
There are two such chosen degree-one primes above each rational
\(p\equiv1\pmod3\). Define \(D=|\PP(X)|\).
For a denominator prime \(\sigma\), the symbol
\(\chi_\sigma(u)=(u/\sigma)_6\) is the sextic residue symbol,
and is zero whenever \(\sigma\mid u\).
The residue matrix has \emph{numerator row} and \emph{denominator
column}:
\[
 (\rX)_{\rho,\sigma}=\chi_\sigma(\rho),\qquad
 (\rX)_{\rho,\rho}=0.
\]

Write \(e(z)=e^{2\pi iz}\), and for \(p=\Nm(\sigma)\) define
the rational-normalized Gauss sum
\[
 \gtilde(\sigma)=p^{-1/2}\sum_{a=1}^{p-1}
       \chi_\sigma(a)e(a/p).
\]
The rational additive character is essential; it differs from the
trace character used in some composite-modulus comparisons.

A twist class \(\kappa=\prod_\pi\pi^{e_\pi}\) has finitely
many distinct factors, \(1\le e_\pi\le5\), taken modulo sixth powers.
The factors may be split or inert, subject to the original
coprimality convention. Put
\[
 n=\Nm_6(\kappa)=\prod_\pi\Nm(\pi)^{e_\pi},\qquad
 w_\kappa(\sigma)=\overline{\gtilde(\sigma)}\chi_\sigma(\kappa).
\]
Define, where denominators are nonzero,
\begin{align}
 r_\kappa(X)&=\frac{\|\rX w_\kappa\|}
                  {\sqrt D\,\|w_\kappa\|},&
 \epsilon_\kappa(X)&=\frac{r_\kappa(X)^2-1}{r_1(X)^2-1},
 \label{eq:r}\\
 s_\kappa(X)&=n^{1/6}\epsilon_\kappa(X),&
 t&=\frac nX.\label{eq:st}
\end{align}
A value \(r=1\) is a natural random-phase benchmark, but a finite
comparison must use the symmetry-matched random control, not an
unrestricted independent random matrix.

\fig{primary_lattice}{.80}{Lattice plot of sample primary Eisenstein
prime representatives in the RS-S019 class list, from
\texttt{results/rs\_s019\_moduli\_2026-10-26.json}
(\texttt{30965467}). This is illustrative of indexing, not
a density measurement.}{lattice}

\subsection{Experiment design, controls and computational validation}
The programme used committed pre-data predictions with explicit
tolerances. Discovery, new-class holdouts and new-scale holdouts
are distinct. Failed tests were recorded without retrospectively
changing their criteria. In particular, the RS-S015
norm-only proposal and RS-S017 per-class fifth-power test failed;
neither is described as a success merely because a weaker
grouped statistic was subsequently supported.

The null has the same matrix dimension and reciprocity/conjugation
structure as the observed ensemble. Native finite-field and
matrix-free implementations were cross-checked against existing
engine outputs; the RS-S021 reproducibility file records a
relative difference of at most \(5.27\times10^{-15}\)
at its specified validation scale.
% Source: reciprocity_spectra/results/rs_s021_large_x_2026-10-26.json, validation.X320000.max_relative_r_difference.
This is an agreement of two numerical implementations, not a
theorem about an infinite family.

At large \(X\), a matrix-free native engine evaluates matrix--vector
products without materializing the full dense \(D\times D\) array.
Its code, run parameters, platform and control seeds are preserved
with the result JSON. We do not treat repeated views of the same
data as independent replications.

\fig{prereg_timeline}{.96}{Chronology of finite prediction outcomes:
RS-S015 failure; RS-S016 exploratory regime study; RS-S017 mixed
pass/fail; RS-S019, S020 and S021 registered passes of their specified
decisions. Source score blocks in the six corresponding
\texttt{results/rs\_s0**.json} files at \texttt{30965467}.
Timeline is schematic, not a new experiment.}{timeline}

\begin{table}[htbp]\centering\small
\begin{tabular}{@{}llll@{}}\toprule
Session & Frozen test & Outcome & Preserved caveat\\\midrule
S015 & PROV-RS14 P1--P3 & FAIL/PASS/FAIL & precursor refuted\\
S017 & Q1,Q2,Q4 & PASS & fresh classes\\
S017 & Q3 & FAIL & per-class \(r<1\) clause\\
S019 & prime/all-moduli & PASS & finite comparison\\
S020 & H-minus & PASS & grouped fifth classes\\
S021 & new-scale A, B & PASS & finite \(X\le4\cdot10^6\)\\
\bottomrule\end{tabular}
\caption{Selected frozen prediction outcomes. The complete
criteria and provenance reside in
\texttt{research/CONJECTURE\_REGISTER.md} and
\texttt{reciprocity\_spectra/results/} (\texttt{30965467}).}
\label{tab:pred}
\end{table}

\section{Finite numerical results}\label{sec:results}
\subsection{Baseline and collapse}
The baseline \(r_1\) increases over the measured sequence of cutoffs.
An increase in the baseline is \emph{not} assumed at unobserved
scales and does not prove a positive limiting gap.
\fig{r1_growth}{.85}{Recorded baseline strengths over cutoffs
\(1.6\times10^5\) to \(4\times10^6\). Sources: RS-S016
\texttt{analysis.r\_u1} and RS-S021 \texttt{score.r\_u1},
\texttt{results/rs\_s016\_x\_dependence\_2026-10-26.json}
and \texttt{rs\_s021\_large\_x\_2026-10-26.json}
(\texttt{30965467}).}{baseline}

The exponent-class curves plotted against \(\log_{10}(n/X)\)
show a marked change between the small-aspect enhancement
regime and the large-aspect region. The frozen reference
\(F\) used for subsequent comparisons is an interpolated
discovery bin table from RS-S016, not a derived analytic formula.
\fig{exponent_profiles}{.91}{Exponent-resolved bin means from
RS-S016's \texttt{analysis.collapse\_bins\_t\_eq\_N6\_over\_X};
unlike the held-out points in Figure~\ref{fig:collapse},
these bins are exploratory discovery summaries.
Source: \texttt{rs\_s016\_x\_dependence\_2026-10-26.json}
(\texttt{30965467}).}{profiles}

\subsection{Fresh held-out classes and new scales}
The RS-S017 tests used fresh twist classes rather than only those
that suggested the collapse. The recorded Q1, Q2 and Q4 criteria
passed while Q3 failed: \(30\) of \(224\) class--scale points
had \(r\ge1\), contrary to its originally universal
per-class clause. This failure remains part of the record.
% Source: results/rs_s017_fresh_classes_2026-10-26.json score.Q3.r_failures, n_with_5 (two scales), and score.Q1/Q2/Q4.

At \(X=2\times10^6\) and \(4\times10^6\), the
matrix-free RS-S021 runs found all \(12\) scored bins within
their registered \(\pm0.05\) difference from the frozen
profile. This means the \emph{pre-registered finite test}
passed; it cannot prove uniform convergence in \(X\).
% Source: results/rs_s021_large_x_2026-10-26.json score.A.bins and score.A.outcome.
\fig{holdout_residuals}{.88}{RS-S021 holdout mean
minus frozen-profile mean, with the precommitted
\(\pm0.05\) threshold. Marker areas represent bin sample
counts. Source: \texttt{rs\_s021\_large\_x\_2026-10-26.json},
\texttt{score.A.bins} (\texttt{30965467}).}{residuals}

\subsection{Exponent five and the prime-modulus control}
Fifth-power twists have a negative large-\(t\) excess,
with the registered RS-S021 grouped means
\(-1.0581\) and \(-1.0469\) at the two new cutoffs.
% Source: results/rs_s021_large_x_2026-10-26.json score.B.mean_s_pi5_t_ge_3.
RS-S020 independently found \(-0.9188\) with standard
error \(0.0828\) on \(88\) fresh prime classes
with \(400<\Nm(\pi)\le1000\), consistent with but not
establishing a limiting value \(-1\).
% Source: results/rs_s020_h_minus_2026-10-26.json score.mean_s, score.se, score.n, predictions.
\fig{exponent_five}{.90}{Individual fifth-power classes at
the two RS-S021 new scales. The horizontal reference
\(-1\) is a conjectured large-aspect value, not a fit or
an established limit. Source:
\texttt{rs\_s021\_large\_x\_2026-10-26.json},
\texttt{runs[].classes} (\texttt{30965467}).}{pi5}

Comparing prime denominators to all squarefree denominators
in the trace-normalised RS-S019 variant, the mean fifth-power
excess was \(-1.0812\) versus \(-1.2224\) for the registered
groups. The frozen test classified the deficit as intrinsic
rather than a prime-sieve artefact. This does not equate the
trace and rational normalisations or prove an analytic mechanism.
% Source: results/rs_s019_moduli_2026-10-26.json score.prime.pi5_mean_s and score.unrestricted.pi5_mean_s, score.P.outcome.
\fig{moduli}{.92}{Profiles for prime and all-squarefree denominator
sets at \(X=2\times10^5\). These are RS-S019 \emph{trace}
comparison runs, not the frozen rational-prime matrix.
Source: \texttt{rs\_s019\_moduli\_2026-10-26.json},
\texttt{runs[0].prime/unrestricted.classes} (\texttt{30965467}).}{moduli}

\subsection{Parity structure and the small-\(t\) edge}
The row-prime partition by \(j\) with
\(\chi_\pi(\rho)=\zeta_6^j\), especially for exponent \(2\),
shows an exploratory dependence on
\(\Nm(\pi)\bmod12\). Aggregation by row partition
must not be read as a proven Hecke coefficient formula.
\fig{rho_parity}{.86}{Exploratory means across the row partitions
for exponent \(2\), separated by prime norm modulo \(12\),
in the \emph{rational} setting. Source:
\texttt{rs\_s020\_rho\_partition\_2026-10-26.json},
\texttt{settings.rational.classes} (\texttt{30965467}).
No error bars or inferential claim are assigned.}{parity}

The RS-S021 small-aspect aggregate \(K_1\) rose from
\(0.9612\) at \(X=10^6\) to \(0.9701\) and \(0.9747\)
at the two larger cutoffs. The finite test was classified
as RISING, but \(F(0^+)=1\) is left as an open question.
% Source: results/rs_s021_large_x_2026-10-26.json score.H+, stored_trend_K1.
\fig{plateau}{.85}{Recorded low-aspect band mean
at the six measured cutoffs, compared to the conjectural
boundary value \(1\). Source: RS-S021
\texttt{stored\_trend\_K1}, \texttt{score.H+.K1}
(\texttt{30965467}).}{plateau}

\subsection{Normalisation audit}
A control-referenced definition and the frozen
\(r_1\)-normalised excess show similar scale collapse.
Without the \(r_1\) normalisation, the median across-\(X\)
spread grows to \(0.0671\), versus \(0.0219\) in the
frozen statistic. This is a finite artefact check, not
an asymptotic universality theorem.
% Source: results/rs_s022_cand002_audit_2026-10-26.json normalizations.*.median_across_X_spread_of_bin_means.
\fig{normalisation}{.86}{Across-cutoff bin spread in the three
RS-S022 definitions (frozen, control-referenced, raw). Source:
\texttt{rs\_s022\_cand002\_audit\_2026-10-26.json},
\texttt{normalizations} (\texttt{30965467}).}{normalisation}

\section{Frozen conjectures and boundary distinctions}\label{sec:conj}
The following restates CAND-002 from
\texttt{research/CONJECTURE\_REGISTER.md}, including the
2026-10-27 correction that its third statement is \emph{standalone}.
The original frozen wording and historical clarifications
are preserved in the repository. No statement in this
section is claimed as a theorem.

\begin{conjecture}[CAND-002(1), compact-aspect crossover]\label{conj:one}
There exists \(F:(0,\infty)\to\mathbb R\),
independent of the twist class \(\kappa\), such that for
bounded-factor classes whose exponents all belong to
\(\{1,2,3,4\}\),
\[
 s_\kappa(X)-F\!\left(\frac{\Nm_6(\kappa)}X\right)\longrightarrow0
\]
as \(X\to\infty\), uniformly when \(\Nm_6(\kappa)/X\)
ranges over any fixed compact subset of \((0,\infty)\).
Moreover \(\liminf_{t\to0^+}F(t)>0\), \(F\) is decreasing
on \([10^{-2},1]\), and
\(\limsup_{t\to\infty}F(t)\le0\).
\end{conjecture}

\begin{conjecture}[CAND-002(2), fifth-power large aspect]\label{conj:two}
For \(\kappa=\pi^5\), for every \(\delta>0\)
there exist \(N_0\) and \(T\) such that
\[
 |s_{\pi^5}(X)+1|<\delta
\]
whenever \(\Nm(\pi)\ge N_0\),
\(\Nm(\pi)^5/X\ge T\), and \(X\) is sufficiently large.
\end{conjecture}

\begin{conjecture}[CAND-002(3), standalone fixed classes]\label{conj:three}
For every fixed admissible class \(\kappa\) with
no exponent \(5\),
\(\liminf_{X\to\infty}r_\kappa(X)>1\).
For every fixed prime \(\pi\),
\(\limsup_{X\to\infty}r_{\pi^5}(X)\le1\).
\end{conjecture}

The registry's historical version of (3) phrased it as
a corollary conditional on a sextic Patterson-type
first-moment hypothesis \(H_P\).
The later logical-bridge audit showed that the
\emph{compact-aspect} information in (1) and the
\emph{large-aspect} information in (2) do not control
fixed classes, where \(t\to0\). The 2026-10-27
restatement is the authoritative standalone
fixed-class conjecture. Even corrected \(H_P\) has
errors too large on the typical row scale; see
Section~\ref{sec:heuristic}.
\fig{regime_map}{.90}{The fixed-class boundary is reached by
holding \(\Nm_6\) constant as \(X\) increases, whereas
Conjecture~\ref{conj:one} is uniform on compact positive
\(t\) intervals and Conjecture~\ref{conj:two} concerns
large \(t\). This is a \emph{conceptual} parameter diagram,
not measured data.}{regimes}

Open empirical questions from the frozen record are
(a) whether \(F(0^+)=1\);
(b) the small-aspect fifth-power behavior;
(c) the prime-specific unit-class selection; and
(d) the exponent-two row-parity concentration.
All four are questions, not predictions scored in this paper.

\section{Exact finite core and paper-level identities}\label{sec:proof}
\subsection{Lean-verified finite algebra}
The pinned Lean~4 / Mathlib project under
\texttt{formalisation/} contains exact finite theorems,
not asymptotic estimates. The twenty-six independently
compared selections are listed with precise names and
hypotheses in \texttt{formalisation/README.md};
the comparison includes:
\texttt{sexticExp\_mul\_sixth\_pow},
\texttt{sexticExp\_pow\_five},
\texttt{rawGauss\_modulus\_sq},
\texttt{shiftedGauss\_e4}, and
\texttt{residueMatVec\_frozen\_orientation}.
The checked source baseline is commit
\texttt{8b050b06d1195491b17b92e68b2d92d9d2478cb0}
and the independent Comparator/NanoDa success run is
\href{https://github.com/jfairfaxball-348/PRIMES/actions/runs/40569785577}{40569785577};
the corresponding full scoped cloud workflow is
\href{https://github.com/jfairfaxball-348/PRIMES/actions/runs/40569785692}{40569785692}.
The code also reports \(1764/1764\) agreement in its
finite exponent-faithfulness fixture.
% Source: PROGRAM_STATUS.md RS-S040 checkpoint; formalisation/README.md (main 30965467).
The formal residue convention maps
\(\omega\mapsto r=-a/b\) in \(\mathrm{ZMod}(p)\),
with explicit prime/root/divisibility hypotheses;
it is \emph{not} a Mathlib construction of
the quotient \(\OO/(\sigma)\).
The frozen CAND-002 statements are defined as
Lean propositions, with \emph{no proofs} of
their asymptotic parts. Palomar registration is
not asserted.

\subsection{Exact shifted row and signed Gram identities}
Put \(G_\sigma=\gtilde(\sigma)\),
\(q_\kappa(\sigma)=\overline{G_\sigma}\chi_\sigma(\kappa)\),
and \(U_\kappa=\{\sigma\in\PP(X):\chi_\sigma(\kappa)\ne0\}\);
let \(D_\kappa=|U_\kappa|\).
For the positive-sign rational additive character,
\[
 g_\sigma(u)=\sum_{a\bmod\sigma}
    \chi_\sigma(a)e(\operatorname{res}_\sigma(ua)/\Nm\sigma).
\]

\begin{proposition}[Finite row identity, exact]\label{prop:row}
For every indexed row \(\rho\), with the zero convention,
\[
 (\rX w_\kappa)_\rho=
 \overline{\sum_{\sigma\in\PP(X)}
              \frac{g_\sigma(\rho\kappa)}{\sqrt{\Nm\sigma}}}.
\]
\end{proposition}
\begin{proof}
For \(\sigma\nmid \rho\kappa\), substitute
\(a\mapsto(\rho\kappa)^{-1}a\) in the Gauss sum:
\(g_\sigma(\rho\kappa)
=\overline{\chi_\sigma(\rho\kappa)}g_\sigma(1)\).
Conjugating and multiplying symbols gives exactly the
\(\sigma\)-summand of the matrix product. If
\(\sigma\mid\rho\kappa\), the shifted sum vanishes by
nontrivial multiplicative-character orthogonality,
and so does the corresponding matrix term.
This includes the zero diagonal.
\end{proof}

Define a Hermitian prime-correlation kernel and real signed excess
\begin{align*}
 K_X(\sigma,\tau)
   &=\sum_{\rho\in\PP(X)}
          \chi_\sigma(\rho)\overline{\chi_\tau(\rho)},\\
 O_\kappa
   &=\sum_{\substack{\sigma,\tau\in U_\kappa\\\sigma\ne\tau}}
      q_\kappa(\sigma)\overline{q_\kappa(\tau)}
         K_X(\sigma,\tau).
\end{align*}
\begin{proposition}[Exact zero-diagonal Gram expansion]\label{prop:gram}
For \(D,D_\kappa>0\),
\[
 \|\rX w_\kappa\|^2=D_\kappa(D-1)+O_\kappa,\qquad
 r_\kappa^2=1-\frac1D+\frac{O_\kappa}{DD_\kappa}.
\]
Consequently, wherever \(r_1^2\ne1\),
\[
 s_\kappa=\Nm_6(\kappa)^{1/6}
 \frac{O_\kappa/D_\kappa-1}{O_1/D-1}.
\]
\end{proposition}
\begin{proof}
Expand the squared row sums and interchange finite sums.
For a diagonal column pair \(\sigma=\tau\),
\(|q_\kappa(\sigma)|^2=1\), while the
inner row sum has \(D-1\) unit-modulus terms
and one zero term at \(\rho=\sigma\).
The off-diagonal pairs give \(O_\kappa\).
Conjugate-symmetric pairing makes \(O_\kappa\) real;
normalising gives the formula.
\end{proof}
These results follow RS-S026 and do not assume an
unproved first-moment estimate. Their finite row
counterpart is independently Lean-checked;
the separate Gram expansion is proved here on paper.

\subsection{Reciprocity kernel, sieve and logical limitation}
For distinct columns away from exceptional local primes,
RS-S027 gives an exact reciprocity transform
\[
 K_X(\sigma,\tau)
  =\sum_{\rho\in\PP(X)}
       h_{\sigma,\tau}(\rho)
       \left(\frac{\sigma\tau^5}{\rho}\right)_6,
\]
where \(h_{\sigma,\tau}\) is a \emph{quadratic
local factor modulo \(4\)} determined by the
primary representatives. This is a rearrangement,
not cancellation. The conductor grows with
\(\sigma,\tau\), so the fixed-character prime-ideal
theorem does not establish the required
uniform signed estimate.

The ideal-prime projection in RS-S028 applies
a von Mangoldt weight and separates inert ideals,
prime powers and local exclusions. An exact ideal
Vaughan convolution decomposes the weighted
sum into Type I and Type II pieces.
A separate Dirichlet-series countermodel in
that audit shows that knowledge of a pole
\emph{residue alone} does not determine
a prime-sieved coefficient. These are proofs
of reductions and obstructions, not a proof of
a sextic coefficient asymptotic.

Likewise, the abstract scalar countermodel in
\texttt{RS\_CAND002\_LOGICAL\_BRIDGE\_AUDIT.md}
demonstrates that compact-positive-aspect and
large-aspect hypotheses need not specify the
fixed-class boundary. It is \emph{not}
a counterexample to the arithmetic conjecture.
Detailed mod-\(4\) local signs, trace-to-rational
normalisation corrections, and the ideal
Vaughan terms remain in the referenced
paper-level derivations; no unproved transfer
has been promoted here.

\section{Heuristic interpretation, literature and limitations}\label{sec:heuristic}
\subsection{Proposed mechanism, explicitly heuristic}
Proposition~\ref{prop:row} interprets each matrix row as a
\emph{normalized shifted sextic Gauss-sum first moment}
with shift \(u=\rho\kappa\). It is plausible that
averaging these moments across row primes produces a
crossover when the conductor scale of the shift changes
relative to the prime-modulus cutoff. This is a proposed
mechanism, \emph{not} derived from the Kubota functional
equation or from known theta-coefficient formulas.

The large-\(t\) fifth-power deficit is compatible with
the vanishing condition for some sextic theta
coefficients at fifth prime powers discussed by
Br\"oker--Hoffstein \cite{broeker2013};
compatibility does not determine the \emph{signed}
matrix energy. Similarly, the exploratory exponent-two
row concentration resembles arithmetic local-factor
selection, without constituting a proof.

The scale correction in
\texttt{reciprocity\_spectra/RS\_HP\_SCALE\_ERRATUM.md}
is essential: a Patterson-type sextic prime moment
at size \(X^{2/3}/\log X\) is an asymptotic for
\emph{normalized} sums \(g_\sigma(u)/\sqrt{\Nm\sigma}\).
The analogous raw Gauss-sum scale would be
\(X^{7/6}/\log X\). The former raw-scale
hypothesis in RS-S026/S027 was an error and
its dependent rates must not be invoked as
results for the corrected model.

Even a uniform \(o(X^{2/3}/\log X)\) error in
\(H_P\) exceeds the typical row size
\(\sqrt D\), so it provides no direct control
of the signed excess in
Proposition~\ref{prop:gram}.
A viable proof would require substantially
finer, row-averaged \emph{signed two-point}
estimates in the moving-shift range.
The programme's mathematical proof work on
that question is parked rather than
misrepresented as completed.

\subsection{Relation to existing work}
Residue-symbol matrices and their
characterizations were studied by
Dummit--Dummit--Kisilevsky \cite{ddk2016}.
Patterson, Heath-Brown--Patterson,
Kazhdan--Patterson, and
Chinta--Friedberg--Hoffstein supply part
of the metaplectic/character-sum framework
\cite{patterson1987,hbp1979,kazhdan1984,chinta2003}.
Br\"oker--Hoffstein analyse quartic and
sextic metaplectic coefficients
\cite{broeker2013}. For cubic Gauss sums,
Dunn--Radziwi\l{}\l{} investigate
Patterson's prime-argument bias
\cite{dunn2024}, while David--Dunn--Hamieh--Lin
treat the quartic setting \cite{ddhl2023}.
Blomer--Goldmakher--Louvel provide
large-sieve context \cite{bgl2014};
de Faveri--Dunn--Hoffstein study
nonorthogonality phenomena
\cite{dfdh2026}; Xu--Sang provide recent
norm-level sextic context \cite{xusang2026}.
The analyses here neither claim
independence from these papers nor
infer novelty from a literature search.

\subsection{Limitations and failed antecedents}
The main empirical limitation is that all
arithmetic measurements end at
\(X\le4\cdot10^6\); no data here establish
the function \(F\), uniform limiting
convergence or a fixed-class threshold.
Some plotted points share class lists and
engines; error bars based solely on
independent sampling would be misleading.

The universal norm-only predecessor
CAND-001 is \emph{refuted} by sixth-power
invariance at \(u=64\), and the attempted
repair PROV-RS14 was \emph{refuted} by its
registered RS-S015 class tests.
The per-class fifth-power clause of
RS-S017 Q3 \emph{failed}. The attractive
finite mean deficit is not a justification
for erasing any of these failures.
The H\(_P\) raw-versus-normalized erratum,
the incorrect historical logical-corollary
wording and the absent signed prime-correlation
estimate are further substantive limitations.

\subsection{Availability and responsible-use disclosure}
The code, numerical records, frozen predictions,
session ledger and this manuscript source are in
\url{https://github.com/jfairfaxball-348/PRIMES}.
The primary data are the committed machine-readable
files under \texttt{reciprocity\_spectra/results/};
the plotting program is
\texttt{paper/figures/make\_all.py}. The finite Lean
package is \texttt{formalisation/} on the pinned
toolchain; independently verified selected theorems
are listed in its README. Each plot is reproducible
from the committed JSON, and conceptual diagrams
are labelled as such. No Palomar ID is claimed.

\textbf{AI assistance.}
OpenAI's ChatGPT was used substantially in the
Lean formalisation and assisted with manuscript
drafting. The author is responsible for the
mathematical claims, numerical evidence,
verification, presentation and submission.
The exact underlying model versions were not
consistently recorded. Machine checking of the
finite Lean results does not independently
validate the conjectures or this written text.



\bibliographystyle{amsplain}
\bibliography{references}
\end{document}
